% Speaker notes for opening-slide.svg. Plain LaTeX content for reuse in a notes deck.
\section*{Agora: continue or skip?}

Before the technical walkthrough, ask whether this architecture addresses the
research question closely enough to justify the remaining presentation time.
Continue if Agora is useful as a comparison case; otherwise skip the walkthrough.

\paragraph{Stage placement and group formation are different.}
The authorizer assigns a stage when a worker joins. The inspected client keeps
that assignment during its worker session; we have not established that a
participant can never be reassigned across sessions or by operator action.
Stage affinity is therefore relatively fixed, but averaging membership is not:
each round collects current declarations and independently constructs a roster.
Do not describe this as a permanent fixed group. Agora's inspected matchmaking
does not elect a Hivemind-style leader or use that leader's admission decisions.

\paragraph{Central authority is relevant to research fit.}
The paper explicitly calls the authorizer a single admission and trust authority
and gives it ownership of stage placement. This changes the threat model for
placement and role selection. It does not, by itself, establish that downstream
DHT discovery, roster construction, or parameter aggregation is secure.

\paragraph{Lineage does not establish identical behavior.}
Agora credits Hivemind and reuses its foundations, but its stage-based training,
per-round roster construction, and designated reduction roles need separate
analysis. Frame this as a difference in mechanisms and assumptions rather than
as a criticism that Agora ought to reproduce Hivemind.

\paragraph{Candidate question, not a demonstrated vulnerability.}
Can an admitted participant influence stage placement or obtain a trainer role
outside the intended policy? Stage, trainer/worker role, and reducer duty are
distinct decisions. A contributor worker already performs training computation;
it is not automatically a trainer that chooses data and routes microbatches.
The paper does not document an outside worker becoming a trainer. We have not
verified an exploit or inspected the server-side authorization policy completely.
Our private baseline bypasses authorization, uses one synthetic stage, and does
not test role escalation, stage placement, or a multi-stage trainer pipeline.

\paragraph{Suggested judgment.}
Keep Agora as a comparison case for central placement followed by distributed
coordination. Prioritize it further only if a concrete accessible control-plane
surface fits the study. Role or placement influence is a promising question,
with feasibility still to be established.

\paragraph{Sources.}
Agora paper: \texttt{https://arxiv.org/html/2607.13332v1}, sections 3.3.1--3.3.2
and 3.6.1. Pinned code revision: \texttt{37d29a966528675635161007d7b43388012a98b2}.
See \texttt{agora/src/agora/security/authorizer.py}, method
\texttt{get\_authorization}; and
\texttt{agora\_server/src/agora\_server/core/all\_reduce/matchmaking.py},
method \texttt{look\_for\_group}.

\paragraph{Deck provenance.}
The preserved 15-page input is \texttt{research/agora/Pluraris Agora.pdf}.
The new opening page precedes
the original title slide. Existing page artwork is unchanged. Footer artwork
comes from input page 2, rendered at 1280 pixels wide and cropped to the bottom
48 pixels; it retains the original institutional branding and presentation date.

\section*{Research plan: influence who gets averaged}

This is a proposed research direction. No matchmaking attack, Krum integration,
or successful poisoning experiment is demonstrated by the existing Agora baseline.
The research-plan slide follows the original deck, after the opening fit question.

\paragraph{Hypothesis.}
A modest global attacker population may gain disproportionate representation
among the senders whose vectors actually reach a particular aggregation.
Test whether manipulated discovery or roster formation can increase that local
representation and thereby worsen a robust aggregation rule's behavior.
Stage assignment alone does not establish such influence. The first experimental
gate must demonstrate a working membership manipulation in the private testbed.
Keep reducers honest to isolate this effect from a reducer simply returning lies.

\paragraph{Controls and denominators.}
Hold the global attacker budget, seeds, data, poisoning policy, and network model
fixed between native-membership and manipulated-membership arms. Compare mean
aggregation with a separately specified Krum comparison. Record the global
attacker fraction, each peer's roster, expected senders, actually received
contributions per reducer/coordinate, and the resulting local attacker fraction.
An advertised attacker is not necessarily an effective contributing attacker.
Keep incomplete, failed, and inconsistent rounds visible in the outcome accounting.

\paragraph{Krum's guarantee and demonstrated failure.}
The original Krum paper assumes gradient vectors and establishes resilience under
the count condition $2f+2<n$ together with gradient-distribution and accuracy
conditions. Here $n$ is the aggregation population and $f$ is the tolerated number
of Byzantine inputs. Record Krum's configured tolerance separately from the actual
number of malicious inputs. An actual count above that tolerance, or failure of
the count condition, removes the guarantee; it does not prove an attack succeeds.
Report poisoned-vector selection, output error or direction, held-out loss, and
recovery against matched controls. Selection alone need not imply learning harm.

\paragraph{Integration is an open implementation question.}
The inspected Agora reducer computes a weighted arithmetic mean, not Krum:
\texttt{agora\_server/core/all\_reduce/partition.py},
\texttt{TensorPartReducer.accumulate\_part} and
\texttt{check\_current\_part\_finished}.
Krum requires pairwise distances between complete candidate vectors for the
defined aggregation domain. Applying it independently to reducer chunks is a
different rule and must not be presented as original whole-vector Krum.
First define a controlled whole-vector Krum comparison and its communication cost.
If the domain is only the selected sparse parameter slice, explicitly call this
a Krum-based sparse-state comparison. Agora's asynchronous parameter snapshots
are not the original paper's unbiased gradient estimator setting, so its SGD
theorem does not transfer automatically. A separate gradient-based toy experiment
could test the original setting. Do not claim a drop-in defense or proven bound
for sparse asynchronous state averaging before resolving these choices.

\paragraph{Source.}
Blanchard, El Mhamdi, Guerraoui, and Stainer (2017),
\emph{Machine Learning with Adversaries: Byzantine Tolerant Gradient Descent}.
\texttt{https://papers.nips.cc/paper/2017/file/f4b9ec30ad9f68f89b29639786cb62ef-Paper.pdf}.
Existing native source is pinned to the Agora revision cited above.
